\documentclass[11pt,a4paper]{article}
\usepackage[T1]{fontenc}
\usepackage{lmodern,microtype,xcolor,fancyhdr}
\usepackage[a4paper,margin=25mm]{geometry}
\usepackage{amsmath,amssymb,amsthm,mathtools}
\usepackage[hypertexnames=false,colorlinks=true,linkcolor=blue!45!black,urlcolor=blue!45!black]{hyperref}
\setlength{\emergencystretch}{3em}
\setlength{\parskip}{0.35em}
\pagestyle{fancy}\fancyhf{}
\fancyhead[L]{\small A sharp completeness criterion}\fancyhead[R]{\small Research manuscript}
\fancyfoot[C]{\thepage}\setlength{\headheight}{14pt}
\newtheorem{theorem}{Theorem}
\newtheorem{lemma}{Lemma}
\newtheorem{proposition}{Proposition}
\newtheorem{corollary}{Corollary}
\newtheorem{remark}{Remark}
\newcommand{\diag}{\operatorname{diag}}
\newcommand{\tr}{\operatorname{tr}}
\begin{document}
\title{A Sharp Completeness Criterion for Additive-Difference Metrics\\[4pt]\large A counterfamily, the Dirichlet Fisher metric, and boundary principal curvatures}
\author{Keio Research Notes}
\date{31 October 2026}
\maketitle

\begin{abstract}
For the additive-difference metrics considered in the general framework
of Le Brigant, Preston and Puechmorel, we prove that completeness in every
dimension at least two is equivalent to the divergence of
$\int_1^\infty\sqrt{f'(u)}/f(u)\,du$.
The functions $f_\alpha(u)=u^2/(1+u)^\alpha$, $0<\alpha\le1$, satisfy
all the framework's displayed hypotheses but give complete metrics
exactly at $\alpha=1$. Thus those hypotheses alone do not imply
completeness. The actual Dirichlet Fisher metric is complete; we give
an intrinsic proof with explicit compact bounds for its metric balls.
For its Lorentzian immersion, one beta principal curvature grows as
$x^{-1/2}$ when one parameter $x$ tends to zero with the other fixed,
while the second decays as $x^{1/2}$. This distinguishes the shape
operator from its second-fundamental-form component matrix and corrects
the extrinsic boundary claim used in the original completeness argument.
All conclusions are proved analytically.
\end{abstract}

\paragraph{Scope and source.}
This manuscript concerns the explicit Lorentzian immersion in Alice Le Brigant,
Stephen C. Preston and St\'ephane Puechmorel, \emph{Fisher--Rao geometry of
Dirichlet distributions}, arXiv:2005.05608v2, 19 November 2020,
\url{https://arxiv.org/pdf/2005.05608v2}.
Proposition 8 on PDF page 8 asserts bounded principal curvatures; its proof
specifically says that they tend to zero when the first $n-1$ parameters
tend to zero. Theorem 10 on PDF page 10 uses that assertion to prove
completeness. The same statements occur in the currently downloadable
Preston author copy. The calculations below refute that extrinsic
boundary assertion. They also show that the hypotheses of the source's
general framework do not imply completeness: an explicit one-parameter
family satisfies every displayed hypothesis (3), but is incomplete.
A sharp integral criterion replaces that general conclusion.
The actual Dirichlet Fisher metric is complete, with explicit compact
metric-ball bounds. None of this refutes the negative-sectional-curvature
theorem.

The author copy, arXiv version history, the public Crossref metadata,
and targeted searches for a correction were checked on the date above.
No correction was located in those checks. This is not a claim of
exhaustive bibliographic priority. The publisher's version-of-record
full text was not accessible in this audit.

\section{A sharp completeness criterion}
The source's displayed assumptions (3), on PDF pages 3--4, are
\[
 \begin{gathered}
 f(x)=O(x^2),\quad f'(x)=O(x)\quad(x\downarrow0),\qquad
 f(x)=O(x^2)\quad(x\to\infty),\\
 f''>0,\quad (f/f')''>0.
 \end{gathered}
\tag{A}
\]
Its metric (2) is
\[
 ds^2=\sum_{i=1}^n\frac{dx_i^2}{f(x_i)}
       -\frac{(\sum_i dx_i)^2}{f(\sum_i x_i)},\qquad n\ge2.
\tag{B}
\]
We assume the differentiability required by the displayed expressions,
for example $f\in C^3(0,\infty)$. The conditions at zero and strict
convexity imply continuous endpoint limits $f(0)=f'(0)=0$, $f,f'>0$, and strict superadditivity;
hence (B) is a positive-definite Riemannian metric.

\begin{theorem}[Necessary and sufficient completeness criterion]
Suppose $f$ satisfies (A). Then (B), in any dimension $n\ge2$, is
complete if and only if
\[
 \int_1^\infty\frac{\sqrt{f'(u)}}{f(u)}\,du=\infty.
\tag{D}
\]
The growth assumption at infinity in (A) is not needed for this
criterion; it is included to match the source's stated class.
\end{theorem}
\begin{proof}
Set $G=f/f'$, $t=\sum x_i$ and $B=f(t)-\sum f(x_i)>0$.
The inverse metric is
\[
 g^{ij}=f(x_i)\delta_{ij}+\frac{f(x_i)f(x_j)}B.
\tag{E}
\]
Strict convexity of $f$ gives $0<G(x)<x$. The convex function $G$
therefore extends to $G(0)=0$; $G(x)/x$ is increasing, and $G'>0$.
Let $\ell=\lim_{x\downarrow0}G(x)/x$. Then $\ell\le1/2$. Otherwise choose
$a_0$ with $1/2<a_0<\ell$ and fix a small $r>0$. Since $G(x)\ge a_0x$,
\[
 \log\frac{f(r)}{f(x)}=\int_x^r\frac{du}{G(u)}
 \le\frac1{a_0}\log\frac rx,
\]
and $f(x)\ge f(r)(x/r)^{1/a_0}$ contradicts $f(x)=O(x^2)$.

Choose $r>0$ and $C>0$ so that $G(x)\le3x/4$ and $f(x)\le Cx^2$
for $0<x\le2r$. For $0<m\le r$,
\[
 \frac{f(2m)}{f(m)}=
 \exp\left(\int_m^{2m}\frac{du}{G(u)}\right)\ge2^{4/3}.
\]
The defect is increasing in every coordinate because
$\partial_{x_i}B=f'(t)-f'(x_i)>0$. Hence, if $m=\min x_i$,
\[
 B\ge f(2m)-2f(m)\ge(2^{4/3}-2)f(m).
\]
At an index attaining the minimum, (E) consequently gives
$g^{ii}\le C_0m^2$ for $m\le r$, where
$C_0=C[1+(2^{4/3}-2)^{-1}]$.
The clipped function $\log[r/\min(m,r)]$ thus has metric slope at most
$\sqrt{C_0}$. Here $m$ is locally Lipschitz along any piecewise $C^1$
curve, and almost everywhere its derivative is the derivative of a
coordinate attaining the minimum. Fix a starting point $a$ and a
radius $R\ge0$. On every curve from $a$ of length at most $R+1$,
all coordinates have the uniform lower bound
\[
 \varepsilon_R=\min\{r,\min_i a_i\}
                  e^{-\sqrt{C_0}(R+1)}>0.
\tag{G}
\]

For any such point and any coordinate $i$, decrease the other
coordinates to one coordinate $\varepsilon$ and zeros. Convexity gives
\[
 B\ge f(x_i+\varepsilon)-f(x_i)-f(\varepsilon)
 \ge\varepsilon f'(x_i)-f(\varepsilon)
 \ge[\varepsilon-G(\varepsilon)]f'(x_i)>0.
\]
Define $H'(x)=\sqrt{f'(x)}/f(x)>0$. Equation (E) yields
\[
 \|dH(x_i)\|_{g^{-1}}^2
 =\frac{f'(x_i)}{f(x_i)}+\frac{f'(x_i)}B
 \le\frac1{G(\varepsilon)}
        +\frac1{\varepsilon-G(\varepsilon)}.
\tag{F}
\]
If (D) holds, $H(x)\to\infty$. Use $\varepsilon=\varepsilon_R$ in (F)
and let $L_R$ be the square root of its right side. Every point in the
closed ball of radius $R$ admits a curve from $a$ of length at most
$R+1$, and hence satisfies
$H(x_i)\le H(a_i)+L_R(R+1)$ for each $i$. Because $H$ increases without
bound, this gives a finite uniform upper bound for every coordinate.
Together with (G), it places the closed metric ball in a compact
Euclidean box inside the positive quadrant. The positive-definite Riemannian
metric induces the usual manifold topology, so that ball is closed in
the box and compact. Hopf--Rinow proves completeness.

Conversely, suppose the integral in (D) is finite. Put $q=1/f$ and
$w=-q'=f'/f^2$. The relation $G'=1-ff''/(f')^2>0$ gives
\[
 q''=\frac{(f')^2}{f^3}(1+G')>0,
\]
so $w$ is decreasing. Fix the last $n-1$ positive coordinates with sum
$c>0$ and vary $x_1=u$. Along that ray,
\[
 q(u)-q(u+c)=\int_u^{u+c}w(v)\,dv\le c\,w(u).
\]
The ray's length is bounded by $\sqrt c$ times the finite integral in
(D). It escapes to infinity and proves incompleteness.
\end{proof}


\begin{corollary}[An equivalent test without derivatives]
Under the same hypotheses, completeness is also equivalent to
\[
 \int_1^\infty \frac{du}{\sqrt{u f(u)}}=\infty.
\tag{D$'$}
\]
\end{corollary}
\begin{proof}
Convexity and $G(0)=0$ imply $G(1)u\le G(u)<u$ for $u\ge1$.
Since $\sqrt{f'}/f=1/\sqrt{fG}$, the two positive integrands
are bounded above and below by constant multiples of one another.
\end{proof}
For example, whenever the hypotheses hold and $f(u)$ is bounded above
and below by positive constant multiples of $u(\log u)^p$ for large
$u$, completeness is equivalent to $p\le2$. This follows directly
from the integral of $1/[u(\log u)^{p/2}]$; it does not assert that
an arbitrary choice of $f$ with that growth satisfies the hypotheses.

\section{A counterfamily and its complete endpoint}
\begin{theorem}[An explicit counterfamily]
For every $0<\alpha\le1$, the function
\[
 f_\alpha(x)=\frac{x^2}{(1+x)^\alpha}
\]
satisfies (A). For $0<\alpha<1$, metric (B) is incomplete in every
dimension $n\ge2$.
\end{theorem}
\begin{proof}
Direct differentiation gives
\[
 f_\alpha'(x)=\frac{x\{2+(2-\alpha)x\}}{(1+x)^{\alpha+1}},\qquad
 f_\alpha''(x)=
 \frac{2+4(1-\alpha)x+(1-\alpha)(2-\alpha)x^2}
      {(1+x)^{\alpha+2}}>0,
\]
and, with $G_\alpha=f_\alpha/f_\alpha'$,
\[
 G_\alpha(x)=\frac{x(1+x)}{2+(2-\alpha)x},\qquad
 G_\alpha''(x)=\frac{4\alpha}{\{2+(2-\alpha)x\}^3}>0.
\]
The three growth conditions are immediate.
Fix the last $n-1$ positive coordinates, whose sum is $c>0$, and let
$x_1=u\to\infty$. If $q=1/f_\alpha$, the squared speed along this ray is
$q(u)-q(u+c)$. For $u\ge1$,
\[
 -q'(u)=\frac{(1+u)^{\alpha-1}\{(2-\alpha)u+2\}}{u^3}
       \le(4-\alpha)u^{\alpha-3},
\]
so
\[
 q(u)-q(u+c)=\int_u^{u+c}-q'(v)\,dv
       \le(4-\alpha)c\,u^{\alpha-3}.
\]
The ray escapes every compact subset but has total length at most
\[
 \int_1^\infty\sqrt{(4-\alpha)c}\,u^{(\alpha-3)/2}\,du
 =\frac{2\sqrt{(4-\alpha)c}}{1-\alpha}<\infty.
\tag{C}
\]
Its endpoints contain a metric Cauchy sequence with no limit in the
positive quadrant. This proves metric and geodesic incompleteness.
\end{proof}

In particular, $f(x)=x^2/\sqrt{1+x}$ is an elementary fixed
counterexample. For $n=2$ with $x_2=y>0$, (C) becomes the explicit bound
$4\sqrt{7y/2}$ for the escape from $x_1=1$ to infinity.
Thus the completeness assertion for the whole class (A) is false;
this does not imply incompleteness of the specific trigamma metric.

\begin{corollary}[The endpoint of the counterfamily]
Within the family $f_\alpha(x)=x^2/(1+x)^\alpha$, $0<\alpha\le1$,
metric (B) is complete exactly when $\alpha=1$, in every dimension
$n\ge2$.
\end{corollary}
\begin{proof}
The integrand in (D) is asymptotic to
$\sqrt{2-\alpha}\,u^{(\alpha-3)/2}$. Its integral diverges at
$\alpha=1$ and converges for $0<\alpha<1$.
\end{proof}

\section{Elementary properties of the reciprocal trigamma function}
From now on we return to the actual Dirichlet metric. For $x>0$, write
\[
 S_r(x)=\sum_{k=0}^{\infty}(x+k)^{-r},\qquad
 f(x)=\frac1{S_2(x)}=\frac1{\psi^{(1)}(x)},\qquad
 G(x)=\frac{f(x)}{f'(x)}=\frac{S_2(x)}{2S_3(x)}.
\]
Termwise differentiation, justified locally uniformly, gives
\[
 f'(x)=\frac{2S_3(x)}{S_2(x)^2}>0,
 \qquad
 G'(x)=\frac32\frac{S_2(x)S_4(x)}{S_3(x)^2}-1>\frac12.
\tag{1}
\]
The last inequality is strict Cauchy--Schwarz for the sequences
$((x+k)^{-1})_k$ and $((x+k)^{-2})_k$, which are not proportional.
The convergent series at zero gives
\[
 f(x)=x^2+O(x^4),\qquad f'(x)=2x+O(x^3).
\tag{2}
\]
We will also use the elementary bounds
\[
 \frac{x^2}{1+2x^2}<f(x)<x^2,\qquad f(x)<x.
\tag{3}
\]
Indeed, $x^{-2}<S_2(x)<x^{-2}+\zeta(2)<x^{-2}+2$;
and the left-endpoint sum of the strictly decreasing function
$u\mapsto(x+u)^{-2}$ strictly exceeds its integral $1/x$.

\begin{lemma}[Strict convexity from a positive Laplace kernel]
The function $f$ is strictly convex on $(0,\infty)$.
\end{lemma}
\begin{proof}
Put $k(t)=t/(1-e^{-t})$. By Tonelli,
\[
 S_2(x)=\int_0^\infty e^{-xt}k(t)\,dt.
\]
Moreover
\[
 (\log k)''(t)=-\frac1{t^2}
                  +\frac1{4\sinh^2(t/2)}<0\qquad(t>0).
\]
The tilted density $p(t)=e^{-xt}k(t)/S_2(x)$ is therefore strictly
log-concave. Let $U(t)=\int_t^\infty p(s)\,ds$ be its survival function.
For each $a>0$, the ratio $p(t+a)/p(t)$ is strictly decreasing in $t$.
Consequently
\[
 \frac{U(t)}{p(t)}=\int_0^\infty\frac{p(t+a)}{p(t)}\,da
\]
is strictly decreasing, and the hazard $p(t)/U(t)$ is strictly increasing.
Since $U(0)=1$, integration of that hazard gives
$U(t+s)<U(t)U(s)$ for $t,s>0$. Integrating on the positive quadrant,
\[
 \frac12\mathbb E[T^2]
 =\int_0^\infty\!\!\int_0^\infty U(t+s)\,dt\,ds
 <\left(\int_0^\infty U(t)\,dt\right)^2
 =(\mathbb E[T])^2.
\]
All moments are finite. If
$M_j=\int_0^\infty t^je^{-xt}k(t)\,dt$, this is
$M_0M_2<2M_1^2$. Direct differentiation now yields
$f''=(2M_1^2-M_0M_2)/M_0^3>0$.
\end{proof}

It follows that $f$ is strictly superadditive. More explicitly,
$f(0)=0$ and $f'$ is strictly increasing, so
$f(a+b)-f(a)-f(b)>0$ for $a,b>0$.

\section{The two fundamental forms and the actual shape operator}
Let $n\ge2$, $x_i>0$, $t=\sum_i x_i$, and
\[
 B=f(t)-\sum_i f(x_i)>0.
\]
Choose a primitive $\eta'(x)=1/\sqrt{f(x)}$ and the immersion
\[
 \Phi(x)=(\eta(x_1),\ldots,\eta(x_n),\eta(t))
\]
into Minkowski space with metric
$\sum_{i=1}^n du_i^2-du_0^2$. In the graph coordinates $u_i=\eta(x_i)$,
use the tangent basis
\[
 E_i=\partial_{u_i}+\sqrt{\frac{f(x_i)}{f(t)}}\,\partial_{u_0}.
\]
Set $v_i=\sqrt{f(x_i)}$. The first fundamental form in this basis is
\[
 \mathsf g=I-\frac{vv^{\!T}}{f(t)}.
\tag{4}
\]
Reverse the normal used in the source. With
\[
 \nu=-\frac{\sum_i\sqrt{f(x_i)/f(t)}\,\partial_{u_i}+\partial_{u_0}}
                 {\sqrt{1-\sum_i f(x_i)/f(t)}},
\]
the convention $\mathcal S=-d\nu$ gives
\[
 \mathsf{II}=
 \frac1{2\sqrt B}\left(
   \diag(f'(x_i))-
   \frac{f'(t)}{f(t)}vv^{\!T}\right).
\tag{5}
\]
This follows either by differentiating the graph or by reversing the
normal in Proposition 4 of the source. Principal curvatures solve
\[
 \det(\mathsf{II}-\kappa\mathsf g)=0;
\tag{6}
\]
they are the eigenvalues of $\mathsf g^{-1}\mathsf{II}$, or equivalently
of the symmetric matrix $\mathsf g^{-1/2}\mathsf{II}\mathsf g^{-1/2}$.
They are not generally the ordinary eigenvalues of $\mathsf{II}$.

\section{An explicit boundary counterexample and both principal curvatures}
\begin{theorem}[Beta boundary]
Fix $y>0$, put $p=f'(y)>0$, and define
\[
 c(y)=\frac{\sqrt p\,G'(y)}2>0,
 \qquad d(y)=\frac{2G'(y)-1}{4}>0.
\]
At the beta point $(x,y)$, as $x\downarrow0$, the two principal
curvatures with the orientation in (5), ordered with
$\kappa_+\ge\kappa_-$, satisfy
\[
 \kappa_+(x,y)=\frac{c(y)}{\sqrt x}+O(\sqrt x),\qquad
 \kappa_-(x,y)=\frac{d(y)}{c(y)}\sqrt x+O(x^{3/2}).
\tag{7}
\]
In particular the principal curvatures are not globally bounded on the
immersed beta manifold. Reversing the normal changes both signs and
does not affect unboundedness. Nevertheless
$K(x,y)=-\kappa_+\kappa_-\to-d(y)$.
\end{theorem}
\begin{proof}
Use the abbreviations $F=f(y)$, $p=f'(y)$, and $q=f''(y)$.
From (2),
\[
 B=px+O(x^2),\qquad
 \mathsf g_{22}=\frac{p}{F}x+O(x^2),
\]
whereas
\[
 \mathsf{II}_{22}
 =\frac{p^2/F-q}{2\sqrt p}\sqrt x+O(x^{3/2}).
\]
Already the normal curvature in the direction $E_2$ is
\[
 \frac{\mathsf{II}_{22}}{\mathsf g_{22}}
 =\frac{p^2-Fq}{2p^{3/2}\sqrt x}+O(\sqrt x)
 =\frac{c(y)}{\sqrt x}+O(\sqrt x),
\tag{8}
\]
because $G'=1-ff''/(f')^2$. Its positivity follows from (1).
The Rayleigh principle alone therefore proves unboundedness.

For the precise two-eigenvalue result, write
$H=2\sqrt B\,\mathsf{II}$ and use
$\mathsf g^{-1}=I+vv^{\!T}/B$. Direct calculation gives
\[
 \tr(\mathsf g^{-1}H)
 =f'(x)+f'(y)+
 \frac{f(x)f'(x)+f(y)f'(y)-f'(x+y)(f(x)+f(y))}{B}.
\]
Consequently
\[
 \tr(\mathsf g^{-1}\mathsf{II})
 =\frac{c(y)}{\sqrt x}+O(\sqrt x).
\tag{9}
\]
Likewise,
\[
 \begin{aligned}
 \det(\mathsf g^{-1}\mathsf{II})&=
 \frac{f(x+y)f'(x)f'(y)
 -f'(x+y)\{f(x)f'(y)+f(y)f'(x)\}}{4B^2}\\
 &=\frac{p^2-2Fq}{4p^2}+O(x)
 =d(y)+O(x).
 \end{aligned}
\tag{10}
\]
For sufficiently small $x$ the trace and determinant are positive;
the generalized eigenvalues are real since $\mathsf g$ is positive
definite. The two-root formula applied to (9)--(10) gives (7).
The timelike Gauss equation gives $K=-\kappa_+\kappa_-$.
\end{proof}

\paragraph{A fixed explicit ray.}
Take $y=1$ and $x_j=j^{-2}$. Then
\[
 c(1)=\frac12\sqrt{\frac{2\zeta(3)}{\zeta(2)^2}}
 \left(\frac{3\zeta(2)\zeta(4)}{2\zeta(3)^2}-1\right)>0,
 \qquad \kappa_+(j^{-2},1)=c(1)j+O(j^{-1})\to\infty.
\]
The approximate value $c(1)=0.39975353369420685\ldots$ is only an
illustration, not used in the proof. In contrast, every entry of
$\mathsf{II}$ is $O(\sqrt x)$ along this ray. Thus the ordinary
eigenvalues of its component matrix do tend to zero, explaining why
using that matrix as the shape operator gives the opposite conclusion.

\section{A direct intrinsic completeness proof in every dimension}
In the original parameter coordinates the Fisher metric and its inverse
are
\[
 g_{ij}=\frac{\delta_{ij}}{f(x_i)}-\frac1{f(t)},\qquad
 g^{ij}=f(x_i)\delta_{ij}+\frac{f(x_i)f(x_j)}B.
\tag{11}
\]
The inverse follows from Sherman--Morrison and is unrelated to a choice
of ambient normal. Positive definiteness follows directly from
weighted Cauchy--Schwarz:
\[
 g(v,v)\ge\frac{B}{f(t)}\sum_i\frac{v_i^2}{f(x_i)}>0
 \quad\hbox{for }v\ne0.
\]

\begin{lemma}[Two coordinate barriers]
Let $m=\min_i x_i$, $b(u)=f(2u)-2f(u)$, and $n\ge2$. Then
$B\ge b(m)>0$. If $m\le1/4$ and $i$ is an index with $x_i=m$, then
\[
 g^{ii}<\frac52m^2.
\tag{12}
\]
If all $x_i\ge\varepsilon$ with $0<\varepsilon\le1/4$, then for every
index $i$,
\[
 g^{ii}<x_i+\frac{3x_i^2}{2\varepsilon^2},\qquad
 \|d\log(1+x_i)\|_{g^{-1}}^2
 <\frac14+\frac3{2\varepsilon^2}.
\tag{13}
\]
\end{lemma}
\begin{proof}
Since $f'$ increases, each partial derivative of the defect $B$ is
positive. Decreasing all but two coordinates to zero and the remaining
two to $m$ therefore gives $B\ge b(m)$. For $0<u\le1/4$, (3) yields
\[
 b(u)>\frac{4u^2}{1+8u^2}-2u^2\ge\frac23u^2.
\]
At a minimum coordinate, (11) and $f(m)<m^2$ give
$g^{ii}<m^2+m^4/((2/3)m^2)=(5/2)m^2$.
The same defect estimate at $\varepsilon$ and $f(x_i)<x_i$ give the
first inequality of (13). Divide by $(1+x_i)^2$ and use
$x_i/(1+x_i)^2\le1/4$ and $x_i^2/(1+x_i)^2<1$ for the second.
\end{proof}

\begin{theorem}[Quantitative compactness of Fisher metric balls]
Fix a point $a\in(0,\infty)^n$ and $R\ge0$. Define
\[
 \varepsilon_R=\min\{\min_i a_i,1/4\}
                   \exp\!\left(-\sqrt{5/2}\,R\right),
 \qquad
 C_R=\sqrt{\frac14+\frac3{2\varepsilon_R^2}}.
\]
Every point $x$ in the closed Fisher metric ball of radius $R$ about $a$
satisfies
\[
 \varepsilon_R\le x_i\le(1+a_i)e^{C_RR}-1
 \qquad(1\le i\le n).
\tag{14}
\]
Every closed metric ball is compact. Hence the Dirichlet Fisher manifold
is metrically and geodesically complete for every $n\ge2$.
\end{theorem}
\begin{proof}
For any covector $\alpha$ and tangent vector $v$, metric Cauchy--Schwarz
gives
\[
 |\alpha(v)|\le\|\alpha\|_{g^{-1}}\|v\|_g.
\]
Along a piecewise $C^1$ curve, the minimum coordinate $m$ is locally
absolutely continuous and, almost everywhere, its derivative is the
derivative of one of the coordinates attaining that minimum. By (12),
\[
 |d\log m|\le\sqrt{5/2}\,ds\qquad\hbox{where }m\le1/4.
\]
Equivalently, the locally Lipschitz clipped function
$\log[(1/4)/\min(m,1/4)]$ has metric slope at most $\sqrt{5/2}$.
A curve from $a$ of length at most $L$ therefore has $m\ge\varepsilon_L$
at every point. Apply (13) with this fixed $\varepsilon_L$ and integrate:
\[
 \log(1+x_i)\le\log(1+a_i)+C_LL.
\]
This proves (14) for endpoints of curves of length at most $R$.
For distance at most $R$, use curves of length at most $R+\eta$ and let
$\eta\downarrow0$; no prior existence of minimizing geodesics is used.

The box in (14) is a compact subset of the positive quadrant.
The metric is of class $C^3$ and positive definite there, and its topology is
the Euclidean manifold topology. A closed metric ball is a closed subset
of that compact box, and hence compact. Hopf--Rinow completes the proof.
\end{proof}

\paragraph{Logical consequence for the source.}
The counterexample invalidates the uniform extrinsic-curvature argument
used in Proposition 8 and hence that particular proof of Theorem 10.
The intrinsic theorem above establishes the latter conclusion directly.
It does not require a bounded shape operator, any unproved curvature
conjecture, or any finite numerical certificate.
\end{document}
